\begin{table*}[t]
\centering
\caption{Two Reference Lines with the Raw OCR and Each Method's Output on Tesseract Input}
\label{tab:examples}
\footnotesize
\renewcommand{\arraystretch}{1.15}
\begin{tabular}{p{2.3cm}p{7.2cm}p{7.2cm}}
\toprule
 & Example 1 & Example 2 \\
\midrule
Reference & \bn{শাখায় ঝুলিয়া থাকে, রাত্রি কালে আহার নিমিত্ত বহির্গত} & \bn{যুদ্ধেস্থিত কৌরবেরে দেখিয়া অর্জুন। অস্ত্রপাত} \\
Gloss & \textit{``[they] hang from branches, and at night go out for food''} & \textit{``seeing the Kauravas arrayed for battle, Arjuna \ldots''} \\
\midrule
Tesseract (raw) & \bn{\textcolor{red}{শাখার} \textcolor{red}{ঝুলিয়!} থাকে, রাত্রি কালে আহার \textcolor{red}{নিমিত} \textcolor{red}{বহির্ত}} \hfill{\scriptsize 0.11} & \bn{\textcolor{red}{বুদ্ধেছ্বিত} কৌরবেরে \textcolor{red}{দেখিয়।} \textcolor{red}{অর্জ্ন।} \textcolor{red}{অত্্রপাত}} \hfill{\scriptsize 0.16} \\
EasyOCR (raw) & \bn{\textcolor{red}{শখায়} \textcolor{red}{ঝালয়} থাকে, \textcolor{red}{বাত্রি} কালে \textcolor{red}{অহার} নিমিত্ত \textcolor{red}{হয়,}} \hfill{\scriptsize 0.23} & \bn{যুদ্ধেস্থিত কৌরবেরে \textcolor{red}{দেখিয়ী} \textcolor{red}{ভজ} \textcolor{red}{ন} \textcolor{red}{|} অস্ত্রপাত} \hfill{\scriptsize 0.16} \\
\midrule
SymSpell & \bn{\textcolor{red}{শাখার} \textcolor{red}{ঝুলিয়ে!} থাকে, রাত্রি কালে আহার \textcolor{red}{নিহিত} বহির্গত} \hfill{\scriptsize 0.13} & \bn{\textcolor{red}{বুদ্ধেছ্বিত} \textcolor{red}{গৌরবের} \textcolor{red}{দেখিয়ে।} \textcolor{red}{অর্জন।} \textcolor{red}{অত্্রপাত}} \hfill{\scriptsize 0.22} \\
Gemma 2B & \bn{\textcolor{red}{শাখার} \textcolor{red}{ঝুলিয়!} থাকে, রাত্রি কালে আহার \textcolor{red}{নিমিত} \textcolor{red}{বহির্ত}} \hfill{\scriptsize 0.11} & \bn{\textcolor{red}{বুদ্ধেছ্বিত} কৌরবেরে \textcolor{red}{দেখিয়।} \textcolor{red}{অর্জ্ন।} \textcolor{red}{অত্্রপাত}} \hfill{\scriptsize 0.16} \\
Qwen3 1.7B & \bn{\textcolor{red}{শাখার} ঝুলিয়া থাকে, রাত্রি কালে আহার \textcolor{red}{নিমিত} \textcolor{red}{বহির্ত}} \hfill{\scriptsize 0.09} & \bn{\textcolor{red}{বুদ্ধে} \textcolor{red}{ছ্বিত} কৌরবেরে দেখিয়া \textcolor{red}{অর্জ্ঞা} \textcolor{red}{অত্রপাত}} \hfill{\scriptsize 0.20} \\
Phi-4-mini & \bn{\textcolor{red}{শাখার} \textcolor{red}{ঝুলিয়!} থাকে, রাত্রি কালে আহার \textcolor{red}{নিমিত} \textcolor{red}{বহির্ত}} \hfill{\scriptsize 0.11} & \bn{\textcolor{red}{বুদ্ধেছ্বিত} কৌরবেরে \textcolor{red}{দেখিয়।} \textcolor{red}{অর্জ্ন।} \textcolor{red}{অত্্রপাত}} \hfill{\scriptsize 0.16} \\
Llama 3.2 1B & \bn{\textcolor{red}{শাখার} \textcolor{red}{ঝুলিয়!} থাকে, রাত্রি কালে আহার \textcolor{red}{নিমিত} \textcolor{red}{বহির্ত}} \hfill{\scriptsize 0.11} & \bn{\textcolor{red}{বুদ্ধেছ্বিত} কৌরবেরে \textcolor{red}{দেখিয়।} \textcolor{red}{অর্জ্ন।} \textcolor{red}{অত্্রপাত}} \hfill{\scriptsize 0.16} \\
BanglaT5 & \bn{\textcolor{red}{প্রাণীে।।।} \textcolor{red}{এই} \textcolor{red}{জাতিভুক্তি.ের} \textcolor{red}{অন্তর্ভুক্ত।।}} \hfill{\scriptsize 0.81} & \bn{\textcolor{red}{'অন্তস্থল'।} \textcolor{red}{১৯} \textcolor{red}{'।} \textcolor{red}{'অন্ত'।} \textcolor{red}{১৯} \textcolor{red}{'।} \textcolor{red}{'অন্ত'।} \textcolor{red}{১৯}} \hfill{\scriptsize 0.78} \\
TituLLMs 1B & \bn{\textcolor{red}{মাংসানী} \textcolor{red}{বলিয়া} \textcolor{red}{কয়েক} \textcolor{red}{জাতি} \textcolor{red}{আছে।} \textcolor{red}{তাহারা} \textcolor{red}{মাংস} \textcolor{red}{ভক্ষণ}} \hfill{\scriptsize 0.62} & \bn{\textcolor{red}{তুর্থ} \textcolor{red}{বর্ষে} \textcolor{red}{পঞ্চম} \textcolor{red}{সংখ্যা}} \hfill{\scriptsize 0.71} \\
TituLLMs 1B Inst. & \bn{\textcolor{red}{পাঠ্য} \textcolor{red}{তৈরি} \textcolor{red}{করতে} \textcolor{red}{পারে} \textcolor{red}{না।} \textcolor{red}{টেক্সটটি}} \hfill{\scriptsize 0.70} & \bn{\textcolor{red}{যে} \textcolor{red}{কোনো} \textcolor{red}{বিষয়ে} \textcolor{red}{বই,} \textcolor{red}{নিবন্ধ} \textcolor{red}{বা} \textcolor{red}{অন্যান্য} \textcolor{red}{সম্পদ}} \hfill{\scriptsize 0.78} \\
TituLLMs 3B Inst. & \textit{(empty)} \hfill{\scriptsize 1.00} & \bn{\textcolor{red}{একটি} \textcolor{red}{চিত্র।} \textcolor{red}{এটি} \textcolor{red}{শুধুমাত্র}} \hfill{\scriptsize 0.78} \\
\bottomrule
\end{tabular}
\par\smallskip
\begin{minipage}{\linewidth}\footnotesize Pages VT\_961\_k\_0172, 279\_2\_D\_32\_0007. Each row shows the span of the output aligned to the reference line, with words at either end dropped while that lowers its CER; words that differ from the reference are in \textcolor{red}{red}, and the number at the right is the span's CER against the reference. Glosses are approximate.\end{minipage}
\end{table*}
